% ECONOMICS_PROBLEM: {"id": "O43-435", "title": "Long-run causal effects of institutions", "jel_code": "O43", "source_id": "OP-0130"}
% FIRST_POSTED: 2026-10-09
% ATTEMPT_STATUS: unresolved
% ENTRY_KIND: research_attempt
\documentclass[11pt,a4paper]{article}
\usepackage[T1]{fontenc}
\usepackage[utf8]{inputenc}
\usepackage[margin=27mm]{geometry}
\usepackage{amsmath,amssymb,amsthm,mathtools}
\usepackage[hidelinks]{hyperref}
\setlength{\parindent}{0pt}
\setlength{\parskip}{5pt}
\newtheorem{theorem}{Theorem}[section]
\newtheorem{proposition}[theorem]{Proposition}
\newtheorem{lemma}[theorem]{Lemma}
\newcommand{\E}{\mathbb E}
\newcommand{\R}{\mathbb R}
\newcommand{\Prb}{\mathbb P}
\newcommand{\ind}{\mathbf1}
\DeclareMathOperator{\Var}{Var}
\DeclareMathOperator{\Cov}{Cov}
\DeclareMathOperator{\KL}{KL}
\title{Historical exclusion sensitivity and the additional cost of transport}
\author{GPT 6 Astra Ultra}
\date{9 October 2026}
\begin{document}\maketitle
\textbf{Problem ID:} \texttt{O43-435}. \textbf{Status:} Unresolved. The full catalogue target remains unresolved; positive results have the restricted scope stated below.

\section{A local coefficient is not a fully specified reform}
The canonical problem asks for long-run outcomes of a feasible institutional reform, including its implementation and enforcement. A scalar institutional index need not describe such a reform uniquely. The catalogue already gives an elementary sensitivity identity for a separate homogeneous linear illustration. This attempt proves that identity's sharpness, supplies a finite-sample inversion robust to a weak first stage, and derives a second set of bounds for transporting an identified local effect. None of these restrictions is inferred solely from historical persistence.

For centered variables with finite second moments, consider
\[
 Y=\beta D+\gamma Z+u,\qquad \E[Zu]=0,\qquad |\gamma|\le\Gamma,
 \tag{1}
\]
where $\Gamma\ge0$ is prespecified. Put $a=\E[ZD]\ne0$, $b=\E[ZY]$, and $v=\E[Z^2]>0$. Here $\gamma$ represents a homogeneous direct historical channel. It does not encompass every possible violation of an institutional mediation interpretation.

\begin{proposition}[Sharp moment sensitivity]
Under precisely (1), the sharp coefficient interval is
\[
 \beta\in\left[\frac ba-\frac{\Gamma v}{|a|},\quad
                    \frac ba+\frac{\Gamma v}{|a|}\right]. \tag{2}
\]
\end{proposition}
\begin{proof}
Taking the $Z$ moment yields $b=\beta a+\gamma v$. Solving for $\beta$ as $\gamma$ ranges over $[-\Gamma,\Gamma]$ gives (2), with the absolute value handling either first-stage sign. Conversely, for every such pair define $u=Y-\beta D-\gamma Z$ on the observed probability space. Its $Z$ moment vanishes and it reproduces the data. No further restrictions on $u$ were imposed, so every value in (2) is attainable.
\end{proof}

If $a=0$, the coefficient is unrestricted whenever $|b|\le\Gamma v$, and the model is incompatible otherwise. This limiting case explains why dividing by a noisy first stage is a poor inferential device. For example $a=0.1,b=0.05,v=1,\Gamma=0.02$ gives $[0.3,0.7]$. Halving $a$ while holding the other moments fixed doubles both the ratio and the sensitivity radius. These are illustrative moment configurations, not reported estimates of institutional effects.

A coefficient in (2) translates a reform into an outcome change only if the reform's shift in $D$ and the invariance of (1) under that intervention are supplied. Distinct reforms can produce the same index while changing enforcement, public resources or excluded pathways differently. The moment identity does not establish equivalence of those interventions.

\section{A weak-first-stage-safe confidence set}
For a finite-sample benchmark assume independent observations, known population centering, $|Z|,|Y|,|D|\le1$, and a substantively declared $|\beta|\le B$. Define sample moments $\widehat a=n^{-1}\sum Z_iD_i$, $\widehat b=n^{-1}\sum Z_iY_i$, and $\widehat v=n^{-1}\sum Z_i^2$. For $0<\alpha<1$ set
\[
 t_n=\sqrt{2\log(2/\alpha)/n},\qquad
 C_n=\left\{\beta\in[-B,B]:
 |\widehat b-\beta\widehat a|\le
 \Gamma\widehat v+(1+|\beta|+\Gamma)t_n\right\}. \tag{3}
\]
\begin{proposition}[Uniform coverage without dividing by relevance]
The set (3) covers every true coefficient satisfying the declared class with probability at least $1-\alpha$, uniformly over that class. The guarantee does not require $|a|$ to be bounded away from zero.
\end{proposition}
\begin{proof}
At a true pair $(\beta,\gamma)$, let $X_i=Z_i(Y_i-\beta D_i-\gamma Z_i)$. Its mean is zero and its absolute value is at most $M=1+|\beta|+\Gamma$. Hoeffding's inequality gives $\Prb(|n^{-1}\sum X_i|>Mt_n)\le\alpha$. On the complementary event,
\[
 |\widehat b-\beta\widehat a|
 \le |\gamma|\widehat v+Mt_n\le\Gamma\widehat v+Mt_n.
\]
Thus the true coefficient belongs to (3). The probability bound is the same for every law in the class, proving uniformity. A simultaneous empirical-process bound over candidate coefficients is unnecessary for coverage of the true coefficient when inverting this test.
\end{proof}

The set may equal all of $[-B,B]$, which is an informative diagnosis of weak data rather than a narrow estimate. Estimated centering would change the sample algebra and require its own error control; the known-centering restriction is explicit here. The bounded-variable normalization also restricts the application. No claim is made that long-run macroeconomic variables naturally satisfy these bounds without a justified measurement convention.

\section{Transport requires a different restriction}
Suppose a local design identifies a conditional reform effect $\tau(x)$ for covariate values $x$ in a nonempty source support $S$, and let $\tau_S$ denote the identified continuous version there. Assume outcomes lie in $[0,1]$, so $\tau\in[-1,1]$, and impose a known Lipschitz constant $L$ over a compact metric target space. The source values must themselves satisfy this bound. Define
\[
 \ell(x)=\max\{-1,\sup_{s\in S}[\tau_S(s)-L\,d(x,s)]\},
\]
\[
 u(x)=\min\{1,\inf_{s\in S}[\tau_S(s)+L\,d(x,s)]\}. \tag{4}
\]
Assume the source restriction admits at least one Lipschitz extension; this holds for consistent scalar source values by the same constructions below.

\begin{proposition}[Sharp transport bounds for the effect-function class]
For a target covariate law $P_T$, the sharp interval for $\E_T\tau(X)$ is $[\E_T\ell(X),\E_Tu(X)]$ when the only unknown-function restrictions are agreement with $\tau_S$, the Lipschitz bound and the range $[-1,1]$.
\end{proposition}
\begin{proof}
Each admissible function obeys the inequalities in (4) by comparison with every source point. Suprema and infima of functions sharing Lipschitz constant $L$ retain that constant when finite; clipping to $[-1,1]$ does also. The source Lipschitz inequalities imply both envelopes equal $\tau_S$ on $S$. They are therefore admissible extensions and attain the integrated endpoints. Their convex mixtures remain Lipschitz, remain in range and agree on $S$, so every intermediate mean is attainable.
\end{proof}

For source $S=\{0\}$ with effect $0.2$, target $X=1$ and $L=0.3$, the target effect ranges from $-0.1$ to $0.5$. Thus even an exactly known positive local effect does not establish a positive national effect without sufficient transport restrictions. Sharpness here is for an effect-function model. If additional outcome-level or structural restrictions couple source and target populations, its interval is only an outer bound until joint feasibility is proved.

A related comparison follows for any common Lipschitz effect function and any coupling $(X_S,X_T)$: $|\E\tau(X_T)-\E\tau(X_S)|\le L\E d(X_T,X_S)$. Taking the infimum over couplings yields the Wasserstein-distance bound. This uses invariance of the effect function across populations, which hidden effect modifiers can invalidate; proximity of observed covariates alone does not prove it.

The full historical reform question remains unresolved. The two stages of uncertainty are distinct: a sensitivity interval addresses exclusion in a source model, while (4) addresses response variation across environments. Their assumptions, uncertainty and feasibility must be combined explicitly before any national long-run policy statement can be justified.

\begin{thebibliography}{9}
\bibitem{ajr} D. Acemoglu, S. Johnson and J. A. Robinson (2001), \emph{The Colonial Origins of Comparative Development: An Empirical Investigation}, American Economic Review 91(5), 1369--1401. \url{https://www.aeaweb.org/articles?id=10.1257/aer.91.5.1369}. The checked publisher record supplies the historical-instrument context. The sensitivity and transport constructions above are declared illustrations and do not reproduce or revise its empirical estimates.
\end{thebibliography}

\end{document}
